\subsection{Cinemática diferencial e Jacobiano do efetuador (Em desenvolvimento)}

Para relacionar as velocidades das juntas do manipulador robótico com a velocidade espacial\footnote{\,Velocidade dentro do espaço de trabalho do manipulador robótico} do efetuador, utiliza-se o Jacobiano $J(q)$; o qual mapeia as velocidades em cada junta e as transforma em velocidades - lineares e angulares - do efetuador. Considerando
$\vec q$ como o vetor de juntas e
$\dot{q}$ as suas respectivas velocidades, a velocidade espacial do efetuador (linear e angular) no frame da base é dada pela derivada da cinemática direta:

%% ===== EXPLICAÇÃO =====
% O vetor do lado esquerdo (x_ponto) é a velocidade espacial do efetuador
% V_E é a velocidade linear (m/s)
% Omega_E é a velocidade angular (rad/s)

% O vetor q_ponto contém as velocidades das juntas.
% Cada elemento é thata_ponto_i para as juntas revolutas e d_ponto_i para as juntas prismáticas.

% O Jacobiano J(q) é a mtriz 6xn que mapeia velocidades de juntas para velocidades do efetuador.

% Ou seja: sabendo as velocidades das juntas, multiplica-se por J(q) e obtém como o efetuador está se movendo no espaço.

\begin{equation}
\label{eq:jac_map_vel1}
\dot{\vec{x}}_{E} \;=\;
\begin{bmatrix}
\vec{v}_{E} \\
\boldsymbol{\omega}_{E}
\end{bmatrix}
\;=\;
J(q)\,\dot{\vec{q}},
\end{equation}

onde:
\begin{equation}
    \label{eq:jac_map_vel2}
J(q)\in\mathbb{R}^{m\times n}.
\end{equation}

%% ===== EXPLICAÇÃO =====
% Para cada junta "i", a coluna correspondente "J_i" é construida a partir do eixo da junta (z_i-1) e das posições de referencia (p_e, p_i-1).
% p_E = posição do efetuador
% p_i-1 = posição da junta anterior
% z_i-1 = eixo da junta anterior (eixo de rotação para junta revoluta, direção de movimento para junta prismática)
%
% Junta revoluta: o produto vetorial mostra que a velocidade linear do efetuador vem da rotação em torno de z_i-1; a velocidade angular é simplesmente o eixo da rotação.
%
% Junta prismática: não há rotação, portanto omega = 0; o efetuador se move linearmente na direção do eixo da junta.
%
% Portanto, cada coluna do Jacobiano corresponde à contribuição individual de uma junta para a velocidade do efetuador como um todo.
As colunas de $J(q)$, quando são expressas em um mesmo frame (como o $frame \{0\}$), para a $i$-ésima junta são:
\begin{equation}
\label{eq:jac_columns_revoluta}
\text{(junta revoluta)}\;\;
\begin{cases}
J_{v,i} \;=\; {}^{0}\!\vec{z}_{\,i-1} \times \big({}^{0}\!\vec{p}_{\,E} - {}^{0}\!\vec{p}_{\,i-1}\big), \\
J_{\omega,i} \;=\; {}^{0}\!\vec{z}_{\,i-1},
\end{cases}
\end{equation}

E para a junta prismática:
\begin{equation}
    \label{eq:jac_columns_prismatica}
\text{(junta prismática)}\;\;
\begin{cases}
J_{v,i} \;=\; {}^{0}\!\vec{z}_{\,i-1}, \\
J_{\omega,i} \;=\; \vec{0}.
\end{cases}
\end{equation}

%% ===== EXPLICAÇÃO =====
% Qdo o manipulador está preso a uma base não-fixa, a veloc. do efetuador não depende apenas das juntas: há uma contribuição da movimentação da base.

% O termo com a adjunta (Ad) leva em conta a contribuição do mov da base
% Se a base é fixa, omega = 0, sobra apenas o termo J(q)q_ponto.

% A Ad possui os termos: Rotação (R), posição (p) e o produto vetorial (p_x) em forma matricial, matriz anti-simetrica.

% Portanto, esse termo "transporta" a velocidade da base para o frame do efetuador.
Durante a fase operacional, quando o subsistema do manipulador robótico atua como um corpo não-rígido, a torção total no $frame \{0\}$ é resultado da contribuição da base e das juntas via a adjunta $Ad_{{}^{0}T_{B}}$:
\begin{equation}
\label{eq:jac_base_geral}
{}^{0}\!\vec{V}_{E}
\;=\;
\operatorname{Ad}_{\,{}^{0}\!T_E \left({}^{0}\!T_B\right)^{-1}}\; {}^{0}\!\vec{V}_{B}
\;+\;
J_s(q)\,\dot{\vec{q}},
\end{equation}

onde ${}^{0}\!\vec{V}_{B}=\big[\;{}^{0}\!\vec{v}_{B}\;,\;{}^{0}\!\boldsymbol{\omega}_{B}\;\big]^{\mathsf T}$. 

No caso particular ${}^{0}\!\vec{v}_{B}=\vec{0}$:
\begin{equation}
\label{eq:jac_base_particular}
{}^{0}\!\vec{V}_{E}
\;=\;
\operatorname{Ad}_{\,{}^{0}\!T_E \left({}^{0}\!T_B\right)^{-1}}
\begin{bmatrix}
\vec{0}\\[2pt] {}^{0}\!\boldsymbol{\omega}_{B}
\end{bmatrix}
\;+\;
J_s(q)\,\dot{\vec{q}}.
\end{equation}

A adjunta associada a $T=\begin{bmatrix}\mathbf R & \vec p\\ \mathbf 0 & 1\end{bmatrix}\in SE(3)$ é:
\begin{equation}
\label{eq:adjunta_twist}
\operatorname{Ad}_{T}
\;=\;
\begin{bmatrix}
\mathbf R & [\vec p]_\times \mathbf R \\
\mathbf 0 & \mathbf R
\end{bmatrix}.
\end{equation}


%% ===== EXPLICAÇÃO =====
% Forma dual do Jacobiano.
% Se o efetuador aplica uima força (f_E) e um momento (m_E) sobre o ambiente, o Jacobiano transposto mapeia isso para os torques ou forças necessários em cada junta.
%
% Portanto, J leva as velocidades de juntas (veloc do efetuador)
% J^T leva forças no efetuador (torques nas juntas)

\subsection{Forças e torque}
Para mapear as forças e torque das juntas, pode ser expressa matematicamente como:
\begin{equation}
\label{eq:wrench_to_tau}
\boldsymbol{\tau} \;=\; J(q)^{\mathsf T}\,\vec{w}_{E},
\qquad
\vec{w}_{E} \;=\;
\begin{bmatrix}
\vec{f}_{E}\\ \vec{m}_{E}
\end{bmatrix}.
\end{equation}

Se for necessário transportar as forças e torques entre os frames, usa-se a adjunta dual:
\begin{equation}
\label{eq:adjunta_dual1}
\operatorname{Ad}^{*}_{T}
\;=\;
\begin{bmatrix}
\mathbf R & \mathbf 0 \\
[\vec p]_\times \mathbf R & \mathbf R
\end{bmatrix},
\end{equation}

onde:
\begin{equation}
\label{eq:adjunta_dual2}
\vec{w}' \;=\; \operatorname{Ad}^{*}_{T^{-1}}\,\vec{w}.
\end{equation}
%% ===== EXPLICAÇÃO =====
% "m" é a dimensão da tarefga, 6 se for controle completo de posição e orientação.
% Se "rank(J)" cai, o robô perde graus de liberdade efetivos.
% >>> Exemplo: um braço esticado pode perder a capacidade de se mover em certas direções.
% Os comandos podem "explodir" proximo de singularidades.

% O teste numerico é um jeito de estimar numericamente o Jacobiano por diferenças finitas, desloca-se um pouco "q" e ve como x_e varia.
% Isso é util pra detectar singularidades.
\subsection{Singularidades (Em desenvolvimento)}
A configuração $q$ é singular quando $rank J(q) < m$, onde $m$ é a dimensão da tarefa (por exemplo, $m=6$ para controle de torção completa). Uma checagem numérica pode usar diferenças finitas:
\begin{equation}
\label{eq:rank_condition}
\mathrm{rank}\, J(q) \;<\; m,
\end{equation}
com $m$ a dimensão da tarefa (tipicamente $m=6$ para pose completa). Uma aproximação numérica de $J$ pode ser obtida por diferenças finitas:
\begin{equation}
\label{eq:jac_fd}
J(q) \;\approx\;
\frac{\vec{x}_{E}(q+\Delta q) - \vec{x}_{E}(q)}{\Delta q}.
\end{equation}

%% RESUMO GERAL
% Jacobiano = derivada da cinemática direta.
% >>> Conecta velocidades de juntas com velocidades do efetuador.
% Colunas = efeito de cada junta.
% >>> Revoluta: rotação gera veloc linear (via produto vetorial)
% >>> Prismática: só deslocamento linear
% Base móvel = se o manipulador não está fixo, há termos adicionais via adjunta.
% Dualidade
% >>> J(q_ponto): velocidade do efetuador
% >>> J^T(w_e): torques ou juntas necessários
% Singularidades = O Jacobiano perde "posto" e o robô não consegue se mover em todas as direções desejadas.
